% Part: lambda-calculus
% Chapter: lambda-definability
% Section: lambda-definable-recursive

\documentclass[../../../include/open-logic-section]{subfiles}

\begin{document}

\olfileid{lam}{ldf}{ldr}
\sourcecorrection{OLTELAMLDFLDR-001}{ఆంగ్ల మూలంలో ఈ విభాగపు గుర్తింపు \texttt{lam/dfl/ldr}గా ఉంది. ఇదే అధ్యాయంలోని ఇతర విభాగాల గుర్తింపు \texttt{lam/ldf/*}గా ఉండగా, \texttt{dfl}లో అక్షరాల క్రమం తారుమారైంది. అధ్యాయ గుర్తింపుకు అనుగుణంగా \texttt{ldf}గా సరిచేశాం.}
\olsection{\usetoken{S}{lambda definable} ప్రమేయాలు పునరావృత్తమైనవి}

అన్ని పాక్షిక పునరావృత్త ప్రమేయాలు
\tetoken{లాంబ్డాతో నిర్వచించదగినవే}{!!{lambda definable}}
కాదు; దీనికి విలోమమూ నిజమే. అంటే,
\tetoken{లాంబ్డాతో నిర్వచించదగిన}{!!{lambda definable}}
ప్రమేయాలన్నీ పాక్షిక పునరావృత్తమైనవి.

\begin{thm}
  \ollabel{thm:lambda-computable} పాక్షిక ప్రమేయం $f$
  \tetoken{లాంబ్డాతో నిర్వచించదగినదైతే}{!!{lambda definable}},
  అది పాక్షిక పునరావృత్త ప్రమేయం.
\end{thm}

\begin{proof}
  నిరూపణ రూపురేఖలను మాత్రమే ఇస్తాం. ముందుగా $\lambd$-పదాలను
  అంకీకరిస్తాం: శ్రేణులను సంకేతీకరించే సాధారణ ప్రధాన సంఖ్యల
  ఘాతాల పద్ధతితో ప్రతి $\lambd$-పదానికి క్రమబద్ధంగా గోడెల్
  సంఖ్యను కేటాయిస్తాం. తరువాత లాంబ్డా పదపు గోడెల్ సంఖ్య~$t$ను
  ఆర్గ్యుమెంటుగా తీసుకునే పాక్షిక పునరావృత్త ప్రమేయం
  $\fn{normalize}(t)$ను నిర్వచిస్తాం. ఆ పదానికి నియత రూపం
  ఉంటే దాని గోడెల్ సంఖ్యను ఇది ఇస్తుంది; లేకుంటే ఇది
  నిర్వచితం కాదు. అలాగే, సహజ సంఖ్యలకూ వాటికి సంబంధించిన
  చర్చ్ సంఖ్యాంకాల గోడెల్ సంఖ్యలకూ మధ్య మార్పిడి చేసే
  $\fn{toChurch}$, $\fn{fromChurch}$ అనే రెండు పాక్షిక
  పునరావృత్త ప్రమేయాలను నిర్వచిస్తాం.

  ఈ పునరావృత్త ప్రమేయాలను వాడి $f$ను పాక్షిక పునరావృత్త
  ప్రమేయంగా నిర్వచించవచ్చు. $f$ను
  \tetoken{లాంబ్డాతో నిర్వచించే}{!!{lambda define}s}
  $\lambd$-పదం~$F$ ఉందనుకుందాం. $f(n_1,\dots,n_k)$ను
  గణించడానికి ముందుగా $\fn{toChurch}(n_i)$తో సంబంధిత
  చర్చ్ సంఖ్యాంకాల గోడెల్ సంఖ్యలను పొందుతాం. వాటిని
  $\Gn{F}$కు జోడించి $F \num{n_1}\dots\num{n_k}$
  పదపు గోడెల్ సంఖ్యను పొందుతాం. ఇప్పుడు ఆ గోడెల్ సంఖ్యపై
  $\fn{normalize}$ను ప్రయోగిస్తాం. $f(n_1,\dots,n_k)$
  నిర్వచితమైతే, $F \num{n_1}\dots\num{n_k}$కు నియత
  రూపం ఉంటుంది; అది చర్చ్ సంఖ్యాంకమే అయి ఉండాలి.
  నిర్వచితం కాకపోతే ఆ పదానికి నియత రూపం ఉండదు; అందువల్ల
  \[\fn{normalize}(\Gn{F\num{n_1}\dots\num{n_k}})\]
  కూడా నిర్వచితం కాదు. చివరగా, నియత రూపం పొందిన పదపు
  గోడెల్ సంఖ్యపై $\fn{fromChurch}$ను ప్రయోగిస్తాం.
\end{proof}

\end{document}
